% GENERATED FILE. Edit canonical lessons, then run npm run generate:books.
% canonical-source-hash: fnv1a64:ff49cb4059405b86
% canonical-lessons: JA-C249-kokoro, JA-C249-kimochi, JA-C249-seikaku, JA-C249-majime, JA-C249-wagamama, JA-R249-review-248-249, JA-R249-review-final-249

\chapter{Heart and Character}
\label{ch:ja-heart-and-character}
\hlchaptermodality{249}

\begin{chapteropening}
\textbf{By the end of this chapter:} \emph{I can ask about feelings and character.}

Complete the last lesson of chapter 249: I can ask about feelings and character.
\end{chapteropening}

\section[\texorpdfstring{\ja{こころ} (kokoro)}{こころ (kokoro)}]{\ja{こころ} (kokoro) — the heart; the mind}

\label{lesson:JA-C249-kokoro}



\begin{quote}
Before the new one: say the Japanese for a secret, then the Japanese for to doubt.
\end{quote}

\subsection*{You'll want to know: \ja{こころ}}
\textbf{\ja{こころ}} — \emph{kokoro} — \textquotedblleft{}the heart; the mind\textquotedblright{}.

\begin{quote}
\textbf{\ja{そぼは} \ja{こころが} \ja{やさしいですか。}} — \emph{sobo wa kokoro ga yasashii desu ka} — \textquotedblleft{}Is my grandmother kind-hearted?\textquotedblright{}
\end{quote}

\begin{cousinweb}
\textbf{\ja{こころ}} (\emph{kokoro}) is a native Japanese word (wago), inherited rather than borrowed from Chinese; its kanji was attached later as a kun'yomi (JMdict, KANJIDIC2).
\end{cousinweb}

\begin{grammarlens}[title={What you've built}]
One more word for everyday life at A2.
\end{grammarlens}

\subsection*{Your turn}
\begin{itemize}
\raggedright
  \item \emph{Say it:} \emph{kokoro}
  \item \emph{Say it:} \emph{kokoro}, once more
  \item \emph{Say it:} \emph{sobo wa kokoro ga yasashii desu ka} — Is my grandmother kind-hearted?
  \item \emph{Recall:} say the Japanese for a secret, then the Japanese for to doubt, then say \emph{kokoro} again
\end{itemize}

\begin{tcolorbox}[breakable,colback=teal!4,colframe=teal!35!black,title={Before you move on}]
What does \ja{こころ} mean? (\textquotedblleft{}The heart; the mind\textquotedblright{}.) Say it once more.
\end{tcolorbox}
\section[\texorpdfstring{\ja{きもち} (kimochi)}{きもち (kimochi)}]{\ja{きもち} (kimochi) — a feeling; a mood}

\label{lesson:JA-C249-kimochi}



\begin{quote}
Before the new one: say the Japanese for to doubt, then the Japanese for the heart.
\end{quote}

\subsection*{You'll want to know: \ja{きもち}}
\textbf{\ja{きもち}} — \emph{kimochi} — \textquotedblleft{}a feeling; a mood\textquotedblright{}.

\begin{quote}
\textbf{\ja{かのじょの} \ja{きもちが} \ja{わかりますか。}} — \emph{kanojo no kimochi ga wakarimasu ka} — \textquotedblleft{}Do you understand her feelings?\textquotedblright{}
\end{quote}

\begin{cousinweb}
\textbf{\ja{きもち}} (\emph{kimochi}) is a mixed word: Sino-Japanese \emph{ki}, \textquotedblleft{}spirit\textquotedblright{} (Middle Chinese \emph{khj+jH}), borrowed from Chinese, joined to native \emph{mochi}, \textquotedblleft{}holding\textquotedblright{}, from \emph{motsu} (JMdict, KANJIDIC2; Middle Chinese in Baxter's transcription).
\end{cousinweb}

\begin{grammarlens}[title={What you've built}]
One more word for everyday life at A2.
\end{grammarlens}

\subsection*{Your turn}
\begin{itemize}
\raggedright
  \item \emph{Say it:} \emph{kimochi}
  \item \emph{Say it:} \emph{kimochi}, once more
  \item \emph{Say it:} \emph{kanojo no kimochi ga wakarimasu ka} — Do you understand her feelings?
  \item \emph{Recall:} say the Japanese for to doubt, then the Japanese for the heart, then say \emph{kimochi} again
\end{itemize}

\begin{tcolorbox}[breakable,colback=teal!4,colframe=teal!35!black,title={Before you move on}]
What does \ja{きもち} mean? (\textquotedblleft{}A feeling; a mood\textquotedblright{}.) Say it once more.
\end{tcolorbox}
\section[\texorpdfstring{\ja{せいかく} (seikaku)}{せいかく (seikaku)}]{\ja{せいかく} (seikaku) — personality; character}

\label{lesson:JA-C249-seikaku}



\begin{quote}
Before the new one: say the Japanese for the heart, then the Japanese for a feeling.
\end{quote}

\subsection*{You'll want to know: \ja{せいかく}}
\textbf{\ja{せいかく}} — \emph{seikaku} — \textquotedblleft{}personality; character\textquotedblright{}.

\begin{quote}
\textbf{\ja{かれは} \ja{せいかくが} \ja{いいですか。}} — \emph{kare wa seikaku ga ii desu ka} — \textquotedblleft{}Does he have a good personality?\textquotedblright{}
\end{quote}

\begin{cousinweb}
\textbf{\ja{せいかく}} (\emph{seikaku}) is Sino-Japanese (kango): two Chinese characters read with their on'yomi, readings borrowed from Chinese. They are \emph{sei}, \textquotedblleft{}nature\textquotedblright{} (Middle Chinese \emph{sjengH}) and \emph{kaku}, \textquotedblleft{}status, rank\textquotedblright{} (\emph{kaek}) (KANJIDIC2; Middle Chinese in Baxter's transcription).
\end{cousinweb}

\begin{grammarlens}[title={What you've built}]
One more word for everyday life at A2.
\end{grammarlens}

\subsection*{Your turn}
\begin{itemize}
\raggedright
  \item \emph{Say it:} \emph{seikaku}
  \item \emph{Say it:} \emph{seikaku}, once more
  \item \emph{Say it:} \emph{kare wa seikaku ga ii desu ka} — Does he have a good personality?
  \item \emph{Recall:} say the Japanese for the heart, then the Japanese for a feeling, then say \emph{seikaku} again
\end{itemize}

\begin{tcolorbox}[breakable,colback=teal!4,colframe=teal!35!black,title={Before you move on}]
What does \ja{せいかく} mean? (\textquotedblleft{}Personality; character\textquotedblright{}.) Say it once more.
\end{tcolorbox}
\section[\texorpdfstring{\ja{まじめ} (majime)}{まじめ (majime)}]{\ja{まじめ} (majime) — serious; earnest}

\label{lesson:JA-C249-majime}



\begin{quote}
Before the new one: say the Japanese for a feeling, then the Japanese for personality.
\end{quote}

\subsection*{You'll want to know: \ja{まじめ}}
\textbf{\ja{まじめ}} — \emph{majime} — \textquotedblleft{}serious; earnest\textquotedblright{}.

\begin{quote}
\textbf{\ja{こうはいは} \ja{しごとに} \ja{まじめですか。}} — \emph{kōhai wa shigoto ni majime desu ka} — \textquotedblleft{}Is the junior serious about work?\textquotedblright{}
\end{quote}

\begin{cousinweb}
\textbf{\ja{まじめ}} (\emph{majime}) is written with three characters meaning \textquotedblleft{}true\textquotedblright{}, \textquotedblleft{}face\textquotedblright{} and \textquotedblleft{}eye\textquotedblright{}, read as a whole word; its older origin is not settled (JMdict, KANJIDIC2).
\end{cousinweb}

\begin{grammarlens}[title={What you've built}]
One more word for everyday life at A2.
\end{grammarlens}

\subsection*{Your turn}
\begin{itemize}
\raggedright
  \item \emph{Say it:} \emph{majime}
  \item \emph{Say it:} \emph{majime}, once more
  \item \emph{Say it:} \emph{kōhai wa shigoto ni majime desu ka} — Is the junior serious about work?
  \item \emph{Recall:} say the Japanese for a feeling, then the Japanese for personality, then say \emph{majime} again
\end{itemize}

\begin{tcolorbox}[breakable,colback=teal!4,colframe=teal!35!black,title={Before you move on}]
What does \ja{まじめ} mean? (\textquotedblleft{}Serious; earnest\textquotedblright{}.) Say it once more.
\end{tcolorbox}
\section[\texorpdfstring{\ja{わがまま} (wagamama)}{わがまま (wagamama)}]{\ja{わがまま} (wagamama) — selfish; wilful}

\label{lesson:JA-C249-wagamama}



\begin{quote}
Before the new one: say the Japanese for personality, then the Japanese for serious.
\end{quote}

\subsection*{You'll want to know: \ja{わがまま}}
\textbf{\ja{わがまま}} — \emph{wagamama} — \textquotedblleft{}selfish; wilful\textquotedblright{}.

\begin{quote}
\textbf{\ja{あの} \ja{こどもは} \ja{わがままですか。}} — \emph{ano kodomo wa wagamama desu ka} — \textquotedblleft{}Is that child selfish?\textquotedblright{}
\end{quote}

\begin{cousinweb}
\textbf{\ja{わがまま}} (\emph{wagamama}) is a native Japanese compound (wago), inherited rather than borrowed: it is built from \emph{waga}, \textquotedblleft{}one's own\textquotedblright{}, and \emph{mama}, \textquotedblleft{}as it is\textquotedblright{}, and it is usually written in kana (JMdict, KANJIDIC2).
\end{cousinweb}

\begin{grammarlens}[title={What you've built}]
One more word for everyday life at A2.
\end{grammarlens}

\subsection*{Your turn}
\begin{itemize}
\raggedright
  \item \emph{Say it:} \emph{wagamama}
  \item \emph{Say it:} \emph{wagamama}, once more
  \item \emph{Say it:} \emph{ano kodomo wa wagamama desu ka} — Is that child selfish?
  \item \emph{Recall:} say the Japanese for personality, then the Japanese for serious, then say \emph{wagamama} again
\end{itemize}

\begin{tcolorbox}[breakable,colback=teal!4,colframe=teal!35!black,title={Before you move on}]
What does \ja{わがまま} mean? (\textquotedblleft{}Selfish; wilful\textquotedblright{}.) Say it once more.
\end{tcolorbox}
\section[(dialogue)]{People and feelings, a review}

\label{lesson:JA-R249-review-248-249}



\begin{quote}
Say the Japanese for to believe and a lie.
\end{quote}

\subsection*{Your turn}
Say these aloud:

\begin{itemize}
\raggedright
  \item to believe, a lie, a promise, a secret, to doubt
  \item the heart, a feeling, personality, serious, selfish
\end{itemize}

\begin{tcolorbox}[breakable,colback=teal!4,colframe=teal!35!black,title={Before you move on}]
To believe? (\textbf{\ja{しんじる}}, \emph{shinjiru}.) A lie? (\textbf{\ja{うそ}}, \emph{uso}.) A promise? (\textbf{\ja{やくそく}}, \emph{yakusoku}.) A secret? (\textbf{\ja{ひみつ}}, \emph{himitsu}.) To doubt? (\textbf{\ja{うたがう}}, \emph{utagau}.) The heart? (\textbf{\ja{こころ}}, \emph{kokoro}.) A feeling? (\textbf{\ja{きもち}}, \emph{kimochi}.) Personality? (\textbf{\ja{せいかく}}, \emph{seikaku}.) Serious? (\textbf{\ja{まじめ}}, \emph{majime}.) Selfish? (\textbf{\ja{わがまま}}, \emph{wagamama}.)
\end{tcolorbox}
\section[(dialogue)]{Heart and Character — once more}

\label{lesson:JA-R249-review-final-249}



\begin{quote}
Say the Japanese for the heart and a feeling.
\end{quote}

\subsection*{Your turn}
\begin{itemize}
\raggedright
  \item \emph{Say it:} the heart, a feeling, personality, serious, selfish
\end{itemize}

\begin{tcolorbox}[breakable,colback=teal!4,colframe=teal!35!black,title={Before you move on}]
The heart? (\textbf{\ja{こころ}}, \emph{kokoro}.) A feeling? (\textbf{\ja{きもち}}, \emph{kimochi}.) Personality? (\textbf{\ja{せいかく}}, \emph{seikaku}.) Serious? (\textbf{\ja{まじめ}}, \emph{majime}.) Selfish? (\textbf{\ja{わがまま}}, \emph{wagamama}.)
\end{tcolorbox}
